% Results — Securing LLMs Against Prompt Injection
% Academic manuscript section (IEEE / arXiv style)
% Companion to methodology.tex and experimental_setup.tex
% Numbers from frozen Suite A/B and ablation result tables

\section{Results}
\label{sec:results}

We report results on the frozen held-out test partition under the full two-layer protocol, followed by robustness, ablation, and latency measurements.
Throughout, Bypass Rate and True ASR are kept distinct: the former is the fraction of injection prompts left unblocked, whereas the latter additionally requires judged model compliance.
Unless noted otherwise, percentages are computed over 30 test injections and 23 test benign prompts.

\subsection{Main Comparison on the Held-Out Test Partition}
\label{sec:main-results}

Table~\ref{tab:main-results} summarizes the primary comparison.
The unprotected baseline allows every injection through (Bypass Rate \(100\%\)) and yields a True ASR of \(46.67\%\) (14 successful attacks).
Lexical defenses reduce Bypass Rate only modestly and leave True ASR unchanged at \(46.67\%\): the regular-expression filter reaches \(86.67\%\) Bypass Rate, and the keyword heuristic reaches \(96.67\%\).
In other words, on this partition, regex and keyword filtering block a few injections at Layer~1 but do not reduce end-to-end successful attacks relative to the baseline.

\begin{table}[t]
  \centering
  \caption{Primary test-set results (30 injections, 23 benign). Bypass Rate and True ASR are separate Layer~1 and Layer~2 metrics.}
  \label{tab:main-results}
  \begin{tabular}{lrrrr}
    \hline
    Method & Bypass (\%) & True ASR (\%) & FPR (\%) & Successful attacks \\
    \hline
    Baseline & 100.00 & 46.67 & 0.00 & 14 \\
    Regex & 86.67 & 46.67 & 0.00 & 14 \\
    Keyword & 96.67 & 46.67 & 0.00 & 14 \\
    Method~A (context-aware) & 6.67 & 0.00 & 13.04 & 0 \\
    Method~B (learned) & 20.00 & 13.33 & 8.70 & 4 \\
    \hline
  \end{tabular}
\end{table}

Both context-aware methods outperform the lexical baselines by a wide margin.
Method~A---the hybrid handcrafted gate---achieves a Bypass Rate of \(6.67\%\) and a True ASR of \(0\%\) (no successful attacks), at a false-positive rate of \(13.04\%\) (3 of 23 benign prompts blocked).
Method~B achieves a Bypass Rate of \(20.00\%\) and a True ASR of \(13.33\%\) (4 successful attacks), at a lower false-positive rate of \(8.70\%\) (2 of 23).

These figures answer RQ1--RQ3 directly.
Unprotected risk is substantial (True ASR near one half).
Lexical controls are weak on this corpus once end-to-end compliance is measured.
Context-aware sanitization is markedly stronger, with hybrid Method~A providing the best security and soft-only Method~B offering a milder false-positive trade-off.
As shown later in the ablations, Method~A's advantage should be attributed primarily to its hard-trigger layer rather than to ensemble scoring alone.

\paragraph{Bypass versus True ASR.}
The gap between the two metrics is informative.
Under the baseline and lexical defenses, roughly half of all injections become successful attacks, so many sanitizer misses are converted into complied responses.
Under Method~A, the two residual bypasses do not produce judged compliance, so True ASR falls to zero.
Under Method~B, four of the six bypasses become successful attacks.
Reporting only Bypass Rate would therefore overstate residual harm for Method~A and would understate the distinction between filter misses and realized end-to-end success.

\subsection{Usability: Standard and Edge-Case False Positives}
\label{sec:fpr-results}

On the standard benign test partition, regex and keyword filtering incur no false positives, while Method~A and Method~B incur \(13.04\%\) and \(8.70\%\) FPR, respectively (Table~\ref{tab:main-results}).
The more revealing usability probe is the edge-case benign set of 80 technical-but-harmless prompts (Table~\ref{tab:edge-fpr}).

\begin{table}[t]
  \centering
  \caption{Edge-case benign false-positive rates (\(n=80\)).}
  \label{tab:edge-fpr}
  \begin{tabular}{lrr}
    \hline
    Method & Edge FPR (\%) & Blocked / 80 \\
    \hline
    Baseline & 0.0 & 0 \\
    Regex & 2.5 & 2 \\
    Keyword & 0.0 & 0 \\
    Method~A & 36.25 & 29 \\
    Method~B & 7.5 & 6 \\
    \hline
  \end{tabular}
\end{table}

Lexical methods remain nearly inert on this set.
Method~A, however, blocks 29 of 80 edge prompts (\(36.25\%\)), indicating substantial over-triggering on technical language that overlaps superficially with attack vocabulary.
Method~B is far gentler, blocking 6 of 80 (\(7.5\%\)).
Together with the standard-test FPR figures, these results address RQ4: Method~A purchases stronger security at a clear usability cost, whereas Method~B reduces that cost while allowing a nonzero residual True ASR.

\subsection{Paraphrase Robustness}
\label{sec:para-results}

Table~\ref{tab:paraphrase} reports Bypass Rate before and after mild synonym substitution on 50 training injections (sanitizer-only evaluation).
Both context-aware methods remain at \(0\%\) Bypass Rate on the original and paraphrased variants in this diagnostic setting.
Lexical methods degrade slightly: keyword Bypass Rate rises from \(86\%\) to \(90\%\) (\(+4\) percentage points), and regex from \(74\%\) to \(76\%\) (\(+2\) percentage points).

\begin{table}[t]
  \centering
  \caption{Paraphrase robustness (sanitizer-only Bypass Rate on 50 training injections).}
  \label{tab:paraphrase}
  \begin{tabular}{lrrr}
    \hline
    Method & Original Bypass & Paraphrased Bypass & \(\Delta\) (pp) \\
    \hline
    Baseline & 1.00 & 1.00 & 0.00 \\
    Regex & 0.74 & 0.76 & \(+0.02\) \\
    Keyword & 0.86 & 0.90 & \(+0.04\) \\
    Method~A & 0.00 & 0.00 & 0.00 \\
    Method~B & 0.00 & 0.00 & 0.00 \\
    \hline
  \end{tabular}
\end{table}

The paraphrase operator is intentionally mild, so these results should be read as evidence against shallow lexical brittleness rather than as proof of invariance to arbitrary rewriting.
Even under this limited stress, the lexical baselines move in the expected direction, while the context-aware methods do not.

\subsection{Adaptive Rewriting}
\label{sec:adaptive-results}

Adaptive evaluation rewrites each of the 30 test injections into five styles, yielding up to 150 prompts evaluated with the full two-layer pipeline (Table~\ref{tab:adaptive}).
Relative to the static test set, all active defenses degrade, but the ordering is preserved.

\begin{table}[t]
  \centering
  \caption{Adaptive FLAN-T5 rewriting results (up to 150 prompts; full two-layer evaluation).}
  \label{tab:adaptive}
  \begin{tabular}{lrrr}
    \hline
    Method & Bypass (\%) & True ASR (\%) & Successful attacks \\
    \hline
    Baseline & 100.00 & 45.33 & 68 \\
    Regex & 93.33 & 42.00 & 63 \\
    Keyword & 96.67 & 45.33 & 68 \\
    Method~A & 16.67 & 8.67 & 13 \\
    Method~B & 40.00 & 15.33 & 23 \\
    \hline
  \end{tabular}
\end{table}

Under adaptive rewriting, Method~A Bypass Rate rises from \(6.67\%\) on the static test set to \(16.67\%\), and True ASR rises from \(0\%\) to \(8.67\%\) (13 successful attacks).
Method~B Bypass Rate rises from \(20\%\) to \(40\%\), and True ASR from \(13.33\%\) to \(15.33\%\) (23 successful attacks).
Lexical defenses remain close to the unprotected baseline (True ASR \(42.00\%\)--\(45.33\%\)).
Thus Method~A remains the strongest defense under style shift, while Method~B retains a substantial but smaller advantage over regex and keyword filtering.
The increase relative to the static test set confirms that adaptive restatement is a harder condition than the original wording, even when the rewriter is a compact sequence-to-sequence model rather than an unbounded optimizer.

\subsection{Ablation Studies}
\label{sec:ablation-results}

Ablations attribute the main-table gains to specific components on the same held-out test partition.

\subsubsection{Hard-veto ablation for Method A}

Table~\ref{tab:hard-ablation} confirms the hybrid interpretation of Method~A.
The full system achieves \(6.67\%\) Bypass Rate and \(0\%\) True ASR.
Removing all hard vetoes (soft-only Method~A) raises Bypass Rate to \(83.33\%\) and True ASR to \(40.00\%\) (\(+76.66\) percentage points in Bypass Rate), collapsing performance into the same weak regime as soft-weighted scoring alone and showing that the secondary soft layer is insufficient as a primary defense.
Disabling only the semantic hard veto already raises Bypass Rate to \(63.33\%\) and True ASR to \(23.33\%\), identifying that veto as the principal high-confidence component.
Disabling the intent veto or the role-play--keyword veto individually leaves the full-system metrics unchanged on this partition.

\begin{table}[t]
  \centering
  \caption{Hard-veto ablation for Method~A on the held-out test partition.}
  \label{tab:hard-ablation}
  \begin{tabular}{lrrrr}
    \hline
    Variant & Bypass (\%) & True ASR (\%) & FPR (\%) & \(\Delta\) Bypass (pp) \\
    \hline
    Full Method~A & 6.67 & 0.00 & 13.04 & 0.00 \\
    Soft-only (no hard vetoes) & 83.33 & 40.00 & 0.00 & \(+76.66\) \\
    $-$ semantic hard veto & 63.33 & 23.33 & 0.00 & \(+56.66\) \\
    $-$ intent hard veto & 6.67 & 0.00 & 13.04 & 0.00 \\
    $-$ role-play+keyword hard veto & 6.67 & 0.00 & 13.04 & 0.00 \\
    \hline
  \end{tabular}
\end{table}

Two conclusions follow.
First, Method~A should not be presented as a purely ensemble-based scoring system: soft aggregation alone does not deliver its reported security.
Second, among hard vetoes, the semantic trigger is the decisive primary layer under the present test distribution; the other two vetoes contribute comparatively little once that semantic trigger is present.

\subsubsection{Soft-weighted ablation for Method A}

With hard vetoes disabled, leave-one-signal-out soft ablations remain in a weak band (Table~\ref{tab:soft-ablation}).
The full soft-weighted model already allows \(83.33\%\) Bypass Rate and \(40.00\%\) True ASR.
Removing the semantic soft signal produces the largest further degradation (\(+10.00\) percentage points Bypass Rate).
Removing intent or role-play each adds \(6.67\) percentage points.
Removing regular expression, keyword, or perplexity each adds \(3.34\) percentage points.
Removing instruction shift or objective conflict leaves the soft-only totals unchanged on this partition.

\begin{table}[t]
  \centering
  \caption{Soft-weighted ablation for Method~A (hard vetoes off).}
  \label{tab:soft-ablation}
  \begin{tabular}{lrrr}
    \hline
    Variant & Bypass (\%) & True ASR (\%) & \(\Delta\) Bypass (pp) \\
    \hline
    Full soft-weighted & 83.33 & 40.00 & 0.00 \\
    $-$ semantic & 93.33 & 46.67 & \(+10.00\) \\
    $-$ intent & 90.00 & 46.67 & \(+6.67\) \\
    $-$ role-play & 90.00 & 46.67 & \(+6.67\) \\
    $-$ regex / keyword / perplexity & 86.67 & 40.00 & \(+3.34\) \\
    $-$ instruction shift & 83.33 & 40.00 & 0.00 \\
    $-$ objective conflict & 83.33 & 40.00 & 0.00 \\
    \hline
  \end{tabular}
\end{table}

These results reinforce the hybrid claim: soft fusion contributes some secondary signal, especially through semantics, but does not by itself approach the security of full Method~A.

\subsubsection{Feature ablation for Method B}

Table~\ref{tab:learned-ablation} summarizes Method~B feature ablations.
The full learned model matches the Suite~B main result (\(20.00\%\) Bypass Rate, \(13.33\%\) True ASR, \(8.70\%\) FPR).
Semantic features dominate.
Removing all semantic features raises Bypass Rate to \(100\%\) and True ASR to \(46.67\%\), identical to the unprotected baseline on this partition.
Among individual semantic features, removing the top-three mean causes the largest single increase in Bypass Rate (\(+50.00\) percentage points), followed by the maximum similarity (\(+36.67\)) and the benign-contrast term (\(+30.00\)).
Removing intent produces a small increase (\(+3.33\)).
Removing regular expression, keyword, role-play, instruction shift, objective conflict, or perplexity leaves the full-model metrics unchanged on this test partition.

\begin{table}[t]
  \centering
  \caption{Learned feature ablation for Method~B on the held-out test partition.}
  \label{tab:learned-ablation}
  \begin{tabular}{lrrrr}
    \hline
    Variant & Bypass (\%) & True ASR (\%) & FPR (\%) & \(\Delta\) Bypass (pp) \\
    \hline
    Full Method~B & 20.00 & 13.33 & 8.70 & 0.00 \\
    $-$ all semantic & 100.00 & 46.67 & 0.00 & \(+80.00\) \\
    $-$ semantic top-3 & 70.00 & 33.33 & 0.00 & \(+50.00\) \\
    $-$ semantic max & 56.67 & 26.67 & 0.00 & \(+36.67\) \\
    $-$ semantic contrast & 50.00 & 23.33 & 0.00 & \(+30.00\) \\
    $-$ intent & 23.33 & 16.67 & 8.70 & \(+3.33\) \\
    $-$ regex / keyword / role-play / shift / conflict / perplexity & 20.00 & 13.33 & 8.70 & 0.00 \\
    \hline
  \end{tabular}
\end{table}

A zero change on this partition does not prove that a feature is useless in general; it indicates only that the feature did not alter decisions on these 53 test prompts when the remaining features were present.
The positive conclusion is sharper: semantic features are necessary for Method~B's observed gains, and the expanded semantic representation (maximum, top-three, and contrast) carries most of the learned defense.

\subsection{Runtime Overhead}
\label{sec:latency-results}

Table~\ref{tab:latency} reports CPU latency over 50 prompts.
Lexical filters are negligible (keyword mean \(0.02\)~ms; regex mean \(0.22\)~ms).
MiniLM semantic scoring averages \(20.98\)~ms per prompt, and DistilGPT-2 perplexity scoring averages \(45.03\)~ms.
Full Method~A sanitization averages \(64.79\)~ms (median \(46.68\)~ms; 95th percentile \(147.04\)~ms).

\begin{table}[t]
  \centering
  \caption{CPU latency per prompt (\(n=50\)).}
  \label{tab:latency}
  \begin{tabular}{lrrr}
    \hline
    Component & Mean (ms) & Median (ms) & p95 (ms) \\
    \hline
    Keyword & 0.02 & 0.01 & 0.05 \\
    Regex & 0.22 & 0.07 & 0.78 \\
    MiniLM semantic & 20.98 & 17.42 & 35.61 \\
    DistilGPT-2 perplexity & 45.03 & 29.10 & 111.36 \\
    Full Method~A & 64.79 & 46.68 & 147.04 \\
    \hline
  \end{tabular}
\end{table}

Relative to lexical baselines, Method~A therefore adds tens of milliseconds of pre-inference cost on CPU.
This overhead is small compared with typical large-model generation latency, but it is no longer free: the security gains of context-aware sanitization come with a measurable compute premium concentrated in the embedding and perplexity auxiliaries.

\subsection{Summary of Findings}
\label{sec:results-summary}

Across the held-out test partition, robustness suites, and ablations, five empirical findings stand out.
\begin{enumerate}
  \item Lexical defenses alone do not reduce True ASR relative to the unprotected baseline on this corpus.
  \item Method~A achieves the strongest security (True ASR \(0\%\) on the static test set; \(8.67\%\) under adaptive rewriting), but incurs the highest false-positive burden, especially on edge-case benign prompts.
  \item Method~B improves substantially over lexical baselines with a milder usability cost, yet leaves residual successful attacks (\(13.33\%\) True ASR on the static test set).
  \item Method~A behaves as a hybrid defense: hard semantic vetoes provide most protection, while soft weighted scoring is only a complementary layer (soft-only detection falls to roughly \(17\%\), and True ASR rises to \(40\%\)).
  \item Method~B's gains are driven primarily by semantic features; removing them returns performance to the unprotected baseline on this partition.
\end{enumerate}
These findings motivate framing Method~A as a hybrid architecture rather than a pure ensemble scorer.
Scope boundaries and limitations of the evaluation are stated in Section~\ref{sec:limitations}.
